\chapter{The diagnostic experiments}
\label{ch:diagnostics}

\begin{saybox}[title=What to say at the start of this section]
Frame these as elimination, not as vindication. Each experiment tests one
of the explanations that was proposed, and each returns a negative result.
Negative results are what narrow a diagnosis, and the fact that four
independent interventions all fail to move the score is itself the
strongest evidence that the cause lies outside the model.
\end{saybox}

\section{Design of the diagnostic suite}

All diagnostic experiments share a common harness
(\evd{scripts/experiments/overfitting\_diagnosis.py}) so that comparisons
are matched. The shared elements are: the V7 specialist files; the
production exclusion list; a preprocessing pipeline of median imputation
plus standardisation for numeric columns and most-frequent imputation plus
one-hot encoding for categorical columns, \emph{always fitted on the
training partition only}; the seed $42$; and the production context-holdout
split, except where an experiment explicitly changes it.

Six experiments were executed, producing $164$ individual model fits whose
results are stored as CSV files in
\evd{scripts/experiments/overfitting\_diagnosis/}.

\section{Experiment 1 --- reducing model capacity}

\subsection{What was changed}

Four deliberately low-capacity learners replaced the production ensembles,
holding the data and the split fixed:

\begin{itemize}[leftmargin=1.3em, itemsep=1pt]
\item \textbf{Ridge regression} --- linear in the one-hot expanded
  features, $L_2$-penalised, no interactions and no nonlinearity.
\item \textbf{Shallow decision tree} --- a single tree with restricted
  depth and an enforced minimum leaf size.
\item \textbf{Shallow random forest} --- a small, depth-limited forest.
\item \textbf{Small MLP} --- a compact multilayer perceptron.
\end{itemize}

\subsection{A note on the proposal to use logistic regression and CNNs}

Two suggestions in the original feedback require a technical
clarification, which should be made politely but clearly in the meeting.

\emph{Logistic regression} is a classifier: it models
$P(y = k \mid \mathbf{x})$ over a discrete label set. It cannot be applied
to a continuous day-valued target without first discretising that target,
and discretising changes the prediction problem. The honest linear
analogue for a continuous target is regularised linear regression, which is
exactly what ridge, lasso, elastic net, Bayesian ridge and PLS provide, and
all five were run. Logistic regression \emph{was} additionally run, on a
tertile-binned version of the target, and is reported in \Cref{sec:logistic}
as a reformulation rather than as a substitute.

\emph{Convolutional networks} exploit translation-invariant local structure
in grid-like inputs --- pixels, spectra, time series. The input here is a
tabular row of heterogeneous scalars and categorical codes with no spatial
or temporal adjacency, so convolution has no structure to exploit. The
appropriate neural baseline for tabular data is a fully-connected network,
and a small MLP was run in its place. No CNN was implemented, and the
report says so rather than implying one was.

\subsection{Results}

\tbl{tab_exp1_reduced_complexity}{Experiment 1: reduced-capacity models on
the production data and the production split.}

\fig{fig08_complexity_classes}{0.86}{%
  Mean context-holdout test $\rsq$ across the six specialists for every
  model evaluated in the diagnostic suite, spanning a constant-output
  baseline through to the production ensembles.}

\begin{findingbox}[title=Result of Experiment 1]
Capacity reduction does not reduce the score. On soft/general-shelf-life,
ridge regression reaches $\rsq = 0.910$ and a \emph{small MLP} reaches
$\rsq = 0.972$ --- statistically indistinguishable from the production
LightGBM's $0.971$ on the same split. On hard/general, ridge reaches
$0.929$ against the production XGBoost's $0.933$.

A linear model with no interaction terms recovering a score within two
points of a 2000-round gradient ensemble is a strong signal about the
\emph{problem}, not about the models: the mapping from features to target
in this corpus is close to linear-recoverable within the rules present.
\end{findingbox}

Note also that the reduced-capacity models reproduce the safety-endpoint
failures faithfully --- ridge scores $-0.055$ on hard/safety where the
production Random Forest scores $-0.097$. A shared failure across
capacity classes again points at the data rather than at the learner.

\section{Experiment 1b --- the logistic-regression reformulation}
\label{sec:logistic}

To address the logistic-regression suggestion on its own terms, the
continuous target was binned into tertiles per specialist (short / medium /
long shelf life) and a multinomial logistic regression fitted to the
resulting three-class problem.

\tbl{tab_exp3_logistic}{Experiment 1b: multinomial logistic regression on
tertile-binned shelf life, with the empirical cut-points used.}

\fig{fig12_logistic_regression}{0.9}{%
  Logistic regression on the tertile-binned target. High accuracy on the
  quality specialists mirrors the regression result: a linear decision rule
  separates the classes well within the rules present.}

The result mirrors Experiment 1: on soft/general the classifier reaches
$95.2\%$ test accuracy with macro-$F_1 = 0.952$. This is consistent with
the regression finding and inconsistent with the hypothesis that model
complexity generates the apparent performance. It does not, however, test
generalization any more than the regression did --- it is the same split,
so it inherits the same limitation.

\section{Experiment 2 --- changing $K$ in cross-validation}

\subsection{What was changed}

$K$-fold cross-validation replaced the single holdout, at $K = 5$ and
$K = 10$, for ridge regression and LightGBM.

\tbl{tab_exp2_kfold}{Experiment 2: cross-validated $\rsq$ at two values of
$K$, with per-fold values.}

\fig{fig10_kfold_cv}{0.95}{%
  Cross-validated $\rsq$ at $K = 5$ and $K = 10$. The fold-to-fold standard
  deviation is small in every configuration: the folds agree with each
  other because they are structurally equivalent.}

\begin{findingbox}[title=Result of Experiment 2]
Changing $K$ changes essentially nothing: ridge regression on the
aggregated corpus yields $\rsq = 0.9234 \pm 0.0055$ at $K=5$, with per-fold
values $[0.9217, 0.9293, 0.9225, 0.9289, 0.9145]$.

The reason is structural, and it is worth stating explicitly at the
meeting: $K$ controls \emph{how many} partitions the data is cut into, not
\emph{along which variable} it is cut. If every fold still contains rows
from every generation rule, then every fold is the same experiment
repeated. The low standard deviation across folds is not evidence of
stability; it is evidence that the folds are interchangeable.
\end{findingbox}

\section{Experiment 3 --- aggregating all cheese categories}

The three category-specific training sets were concatenated into a single
corpus and one model per algorithm was trained on the union, testing
whether the specialist architecture was responsible for the scores through
narrow within-category homogeneity.

\fig{fig11_granularity}{0.95}{%
  Per-category specialists versus a single aggregated model, same models
  and same protocol.}

\begin{findingbox}[title=Result of Experiment 3]
Aggregation does not remove the effect. Scores remain in the same region,
which is what the generation-rule account predicts: pooling categories
pools their rules too, so the test rows are still drawn from rules the
training set contains.
\end{findingbox}

\section{Experiment 4 --- removing the validation set}

The validation partition was merged into training, giving an $85/15$
train/test split with the test partition left untouched, and every
model--specialist pair was refitted.

\tbl{tab_exp4_validation_removal}{Experiment 4: test $\rsq$ with and
without a separate validation partition, matched pairs.}

\fig{fig09_validation_removal}{0.56}{%
  Each point is one model--specialist pair. Points lie on the diagonal:
  the presence or absence of a validation split does not materially change
  the test score.}

\begin{findingbox}[title=Result of Experiment 4]
The mean change in test $\rsq$ across all matched pairs is $+0.0003$.
Neither the extra training rows nor the removal of the early-stopping
signal has a material effect.

This should be stated carefully: removing a validation set is \emph{not} a
methodological improvement, and this experiment should not be read as
endorsing it. Its value is purely diagnostic --- it eliminates
validation-set design as a candidate explanation.
\end{findingbox}

\section{Experiment 5 --- the full model zoo}

Ten further regressors were added to obtain a complete capacity spectrum:
lasso, elastic net, Bayesian ridge, PLS regression, $k$-nearest neighbours
($k=5$), RBF-kernel SVR, gradient boosting, AdaBoost, extra trees, and a
mean-predicting dummy baseline. With the four from Experiment 1 and the
three production ensembles this gives seventeen configurations.

\tbl{tab_exp3_extra_regressors}{Experiment 5: the extended model zoo under
the original context-holdout protocol, averaged over the six specialists.}

The dummy baseline is the control that makes the rest interpretable: it
scores $\rsq = -0.003$, confirming the metric behaves as expected and that
the other models are doing real work relative to a constant prediction.

\section{Summary of the diagnostic phase}

\begin{table}[htbp]\centering\small
\caption{Diagnostic summary. Four independent interventions on the model
and the protocol; no material change in the score.}
\label{tab:diag_summary}
\begin{tabular}{p{3.3cm}p{4.1cm}p{5.6cm}}
\toprule
Experiment & What changed & Outcome \\
\midrule
1. Reduce capacity & 17 models from linear to 2000-round ensembles &
Scores essentially unchanged; ridge within 2 points of production LightGBM
on soft/general \\
\addlinespace[2pt]
1b. Reformulate as classification & tertile-binned target, multinomial
logistic & $95.2\%$ accuracy on soft/general; same pattern \\
\addlinespace[2pt]
2. Change $K$ & $K = 5$ and $K = 10$ cross-validation &
$\rsq = 0.923 \pm 0.006$; fold variance negligible \\
\addlinespace[2pt]
3. Aggregate categories & one model over all three categories &
No material change \\
\addlinespace[2pt]
4. Remove validation set & $85/15$ instead of $70/15/15$ &
Mean $\Delta \rsq = +0.0003$ \\
\midrule
\multicolumn{3}{p{13.4cm}}{\textbf{Conclusion of the phase:} every
intervention on the model or on the partition \emph{sizes} failed. The
remaining untested variable was the partition \emph{variable} itself.} \\
\bottomrule
\end{tabular}
\end{table}
